\chapter{背景・関連研究}
% 【この章の目的】提案手法を理解するための前提と，研究の位置づけを示す．
% 内容入り例：examples/chapters/02-background-example.tex
% 背景説明と関連研究レビューを混同せず，必要なら節を分ける．

\section{技術的背景}
% 記号，既知の理論，標準的手法を定義する．教科書的説明は必要範囲に絞る．

\section{関連研究}
% 課題設定・仮定・使用情報・評価方法の軸で比較し，本研究との差分を示す．

\section{本研究の位置づけ}
% 既存研究のどの制約を外し，何を引き継ぐかを明記する．
